\documentclass[conference]{IEEEtran}
\usepackage{cite}
\usepackage{amsmath,amssymb,amsfonts}
\usepackage{algorithmic}
\usepackage{graphicx}
\usepackage{textcomp}
\usepackage{xcolor}
\usepackage{hyperref}

\def\BibTeX{{\rm B\kern-.05em{\sc i\kern-.025em b}\kern-.08em
    T\kern-.1667em\lower.7ex\hbox{E}\kern-.125emX}}

\begin{document}
\begin{titlepage}
    \centering
    \vspace*{\fill}

    \includegraphics[width=0.3\textwidth]{universidad-distrital-logo.png} \\[1.5em]
    \textbf{\Large Universidad Distrital Francisco José de Caldas} \\[0.5em]
    \large Faculty of Engineering \\
    Systems Engineering Department \\[3em]

    {\LARGE \textbf{Challenge 3 - Empirical Evaluation of On-Policy vs. Off-Policy Deep Reinforcement Learning in Gravity-Driven Physics Environments: A Comparative Study of DQN and PPO on ALE/Gravitar-v5}} \\[0.5em]
    {\Large \textit{Environment: Gravitar-v5}} \\[4em]

    \textbf{Authors:} \\[0.8em]
    {\large Yader Ibraldo Quiroga Torres} \\
    \texttt{yiquirogat@udistrital.edu.co} \\[0.5em]
    {\large Rosa Alejandra Lopez Lizarazo} \\
    \texttt{ralopezl@udistrital.edu.co} \\[0.5em]
    {\large Diego Alejandro Garzon Rodriguez} \\
    \texttt{dagarzonr@udistrital.edu.co} \\[4em]

    \textbf{Lecturer:} \\[0.5em]
    Carlos Andrés Sierra Virgüez \\[4em]

    \vfill
    Bogotá, D.C., Colombia \\
    March 2026

\end{titlepage}



\maketitle

\begin{abstract}
This research presents a rigorous empirical comparison between two fundamental families of Deep Reinforcement Learning (DRL) algorithms: the value-based Deep Q-Network (DQN) and the actor-critic Proximal Policy Optimization (PPO).
The study centers on the Atari 2600 environment ALE/Gravitar-v5, a platform distinguished by its challenging Newtonian physics, gravitational pull, and the requirement for precise, sustained thrust control.
While the initial Challenge 1 focused on optimizing DQN as a baseline, Challenge 3 introduces PPO to investigate convergence speed, asymptotic performance, and failure modes under a fixed computational budget.
Our results indicate that PPO not only achieves a higher smoothed mean reward (374.7 vs 335.7) but also demonstrates a transformative improvement in wall-clock efficiency, operating at 171 FPS compared to DQN's bottleneck of 6 FPS.
This paper analyzes these findings in the context of on-policy vs. off-policy dynamics and the impact of entropy-regularized exploration in sparse-reward physics simulations.
\end{abstract}

\begin{IEEEkeywords}
Deep Reinforcement Learning, Proximal Policy Optimization, Deep Q-Network, ALE/Gravitar-v5, Computational Efficiency, Physics-based Control.
\end{IEEEkeywords}

\section{Introduction}
The field of Deep Reinforcement Learning (DRL) has witnessed significant breakthroughs since the introduction of the Deep Q-Network (DQN) by Mnih et al., which first demonstrated human-level performance on a wide array of Atari 2600 games.
Despite these successes, certain environments within the Arcade Learning Environment (ALE) remain notoriously difficult for standard agents.
Environments like \textbf{ALE/Gravitar-v5} present a unique set of challenges that extend beyond simple pattern recognition [1], [2].
Unlike games with static platforms or simple collision logic, Gravitar requires the agent to master momentum, inertia, and gravitational fields to navigate between solar systems and destroy reactors [1], [3].
This paper documents the iterative progress of Group 4 in solving the Gravitar environment.
In the first phase (Challenge 1), we implemented a DQN agent, focusing on experience replay and target networks to stabilize the approximation of the optimal action-value function [4].
In the second phase (Challenge 3), we shifted our focus to Proximal Policy Optimization (PPO), an on-policy method designed to achieve the reliability of trust-region methods while being easier to implement and tune [3], [5].
The central research question driving this study is: \textit{``Under a fixed computational budget and on the same environment, does PPO converge faster, reach higher performance, or exhibit different failure modes compared to the DQN agent designed in Challenge 1?
Why?''} [5]. This question is particularly pertinent for Gravitar, where the physics require a delicate balance of ``thrust'' actions.
DQN's $\epsilon$-greedy exploration often struggles with the precise sequences needed for stable flight, whereas PPO's entropy bonus and clipped surrogate objective theoretically offer more stable policy refinements in such sensitive state spaces [6], [7].
The structure of this paper is as follows: Section II provides the theoretical framework for both algorithms;
Section III details the methodology and identical preprocessing used for fair comparison;
Section IV presents the experimental results derived from TensorBoard logs; Section V offers a deep-dive comparative analysis;
and Section VI concludes with insights into algorithmic suitability for high-precision physics tasks.

\section{Theoretical Framework}
\subsection{Deep Q-Networks (DQN)}
DQN represents the off-policy value-based approach, where the agent learns a function $Q(s, a)$ to estimate the expected future rewards for every possible action in a given state.
Key components include the Replay Buffer, which stores past transitions to break temporal correlations, and the Target Network, which provides stable labels for the loss function calculation [4], [8].

\subsection{Proximal Policy Optimization (PPO)}
PPO is an actor-critic method that updates the policy directly while using a critic to estimate the value function.
Its defining feature is the clipped surrogate objective, which prevents the policy from changing too drastically in a single update, thereby avoiding the ``policy collapse'' often seen in standard policy gradient methods [9], [6].
For Group 4, PPO’s implementation must include Generalized Advantage Estimation (GAE) to manage the trade-off between variance and bias in advantage targets [9].
In the context of the Gravitar environment, GAE is particularly important due to the delay between engine thrust actions and the eventual reward or failure (collision).
By setting $\lambda = 0.95$, we balance the high variance of Monte Carlo returns with the bias of temporal difference learning [1], [2].

\section{Experimental Methodology}
To maintain scientific rigor and ensure that performance gains are attributable to algorithmic improvements rather than environment tuning, we adhered to a strict parity protocol across both Challenges.

\subsection{Environment and Task Description}
Group 4 was assigned \textbf{ALE/Gravitar-v5}. In this environment, the agent controls a spacecraft that must navigate through several solar systems.
The primary challenge is the physics engine, which simulates gravity and inertia [3], [4].
The agent must use its limited fuel to counteract gravity and avoid crashing into terrain while destroying enemy reactors.
Unlike Ms. Pac-Man or Phoenix, Gravitar features a multi-objective structure where basic survival requires constant, high-precision action sequences.

\subsection{Preprocessing Pipeline}
Both the DQN and PPO agents utilized a shared preprocessing wrapper to ensure identical observation spaces [2], [5]:
\begin{itemize}
    \item \textbf{Color Reduction:} RGB images were converted to grayscale to reduce the input dimensionality while preserving crucial luminance features.
    \item \textbf{Dimensionality Standardisation:} Original Atari frames were resized to $84 \times 84$ pixels.
    \item \textbf{Temporal Stacking:} A frame-stacking factor of 4 was applied, providing the agent with a short history of movement to infer velocity and acceleration, which is critical for gravity compensation.
    \item \textbf{Action Skipping:} A frame-skip of 4 was implemented to accelerate training and reduce temporal redundancy in the action space [6], [2].
\end{itemize}

\subsection{DQN Experimental Setup (Challenge 1)}
The DQN agent was trained using a standard convolutional architecture [7].
The replay buffer was configured for $100,000$ transitions to handle the memory-intensive nature of Atari observations.
For exploration, a linear epsilon-greedy schedule was used, decaying from $1.0$ to $0.01$ over $50,000$ steps [8], [9].
The target network was updated every $1,000$ steps to provide stable Q-value targets.

\subsection{PPO Experimental Setup (Challenge 3)}
The PPO implementation followed the actor-critic architecture with a shared CNN backbone [10].
The network consisted of three convolutional layers followed by separate actor and critic heads.
Based on the requirements of Challenge 3 [11], the implementation included:
\begin{enumerate}
    \item \textbf{Horizon ($T$):} A rollout length of 1024 steps was selected to capture longer sequences of navigational thrusts [12].
    \item \textbf{Optimization:} $K = 4$ epochs per rollout with a mini-batch size of $128$.
    \item \textbf{Regularization:} An entropy coefficient $c_2 = 0.01$ was applied to prevent premature convergence to a policy that never uses the engine [12].
\end{enumerate}

\subsection{Hyperparameter Sweep and Seed Control}
A systematic sweep was performed for both agents.
For PPO, we evaluated three learning rates ($1 \times 10^{-4}, 3 \times 10^{-4}, \text{ and } 5 \times 10^{-4}$) to determine the optimal step size for Gravitar's sensitive state-action dynamics [13], [14].
All reported experiments were conducted using a fixed seed (Seed 42) to ensure reproducibility, as mandated by the challenge guidelines [6], [15].

\begin{table}[htbp]
\caption{Hyperparameter Comparison: DQN vs. PPO}
\begin{center}
\begin{tabular}{|l|c|c|}
\hline
\textbf{Hyperparameter} & \textbf{DQN (Challenge 1)} & \textbf{PPO (Challenge 3)} \\
\hline
Learning Rate ($\alpha$) & $1 \times 10^{-4}$ & $2.5 \times 10^{-4}$ \\
\hline
Batch Size & $32$ & $128$ \\
\hline
Discount Factor ($\gamma$) & $0.99$ & $0.99$ \\
\hline
Exploration Strategy & $\epsilon$-greedy ($0.01$ final) & Entropy Bonus ($0.01$) \\
\hline
Training Frequency & 4 steps & 1024 steps (rollout) \\
\hline
Target Network/Clipping & 1,000 steps & $\epsilon = 0.2$ \\
\hline
\end{tabular}
\label{tab1}
\end{center}
\end{table}

\section{Experimental Results}
Training was executed until convergence or until the predefined computational budget was reached.
Data was logged via TensorBoard, tracking episodic reward, frames per second (FPS), and training loss.

\subsection{DQN Learning Dynamics}
The DQN agent showed a slow but steady increase in performance.
As shown in the TensorBoard logs [16], the agent initially struggled with the gravity physics, scoring near zero for the first few thousand steps.
Significant learning began once the epsilon-greedy rate decayed below $0.5$.
By step $236,724$, the agent achieved a smoothed mean reward of \textbf{335.7081} [16].

\begin{figure}[htbp]
\centerline{\includegraphics[width=0.48\textwidth]{imagen dqn 1.png}}
\caption{DQN Baseline Metrics: Displays the 335.7 peak reward and the 6 FPS bottleneck.}
\label{fig_dqn1}
\end{figure}

% Page 3 - Continuing Section IV: Experimental Results

\subsection{PPO Learning Dynamics}
The PPO agent exhibited a more stable and ultimately higher-performing learning trajectory in the Gravitar environment.
Initial exploration was aided by the entropy bonus, which encouraged the agent to engage the engine regularly to overcome gravity.
As seen in the PPO logs (Fig. 2), the agent began consistently accumulating points after 150,000 steps.
By the conclusion of the training session at step $500,736$, the PPO agent achieved a smoothed mean reward of \textbf{374.7062}.

\begin{figure}[htbp]
\centerline{\includegraphics[width=0.48\textwidth]{ppo 1.png}}
\caption{PPO Reward and FPS Results: Shows PPO reaching 374.7 reward and maintaining 171 FPS.}
\label{fig_ppo1}
\end{figure}

Internal training metrics provide further insight into PPO's stability. The \textit{explained variance} of the critic reached a smoothed value of $0.406$, indicating that the value function became moderately accurate at predicting returns despite the environment's complexity,.
The \textit{clip fraction} stabilized at approximately $0.336$, suggesting that roughly one-third of the policy updates were restricted to prevent destructive updates, which is vital for survival in Gravitar's collision-heavy maps,.
Furthermore, the \textit{entropy loss} followed a controlled decay, reaching $-1.9571$, which indicates the agent maintained a healthy degree of exploration throughout the run,.

\section{Comparative Analysis}
This section directly addresses the central research question: does PPO outperform DQN in gravity-driven physics environments under a fixed budget?

\subsection{Sample Efficiency and Asymptotic Performance}
Comparing the learning curves, DQN (Challenge 1) initially appeared to score points earlier in terms of environment steps, reaching a peak of \textbf{335.7} near 236,000 steps,.
However, PPO eventually surpassed this threshold, achieving a final reward of \textbf{374.7}.
In terms of sample efficiency, defined as the number of steps required to reach a target score (e.g., 200 points), DQN was marginally faster at the start.
However, DQN's performance plateaued earlier. This suggests that while DQN's off-policy nature allows it to reuse transitions effectively in the early stages, PPO's on-policy updates and clipped objective enable more refined policy improvements that eventually overcome the complex navigation required for higher scores in Gravitar.

\subsection{Computational and Wall-Clock Efficiency}
The most dramatic difference between the two algorithms was observed in throughput, measured in frames per second (FPS).
According to the logged metrics:
\begin{itemize}
    \item \textbf{DQN Throughput:} Stabilized at a smoothed value of only \textbf{6 FPS}.
    \item \textbf{PPO Throughput:} Maintained a robust rate of \textbf{171 FPS}.
\end{itemize}

This represents a \textbf{28.5-fold increase} in computational efficiency for PPO.
The bottleneck in DQN was primarily attributed to the overhead of managing a large replay buffer and the frequent sampling required for off-policy updates on the specific hardware used.
In contrast, PPO’s on-policy nature, which discards transitions after a rollout epoch, allowed it to utilize the CPU/GPU resources far more efficiently.
In real-world time, the PPO agent was able to complete a 500,000-step training session in a fraction of the time required for a 200,000-step DQN session.

\subsection{Training Stability}
Stability was measured by the variance in rewards across the sliding window and the volatility of the loss function.
DQN's training loss was significantly more volatile, fluctuating between 0.002 and 0.01 throughout the run,.
PPO's loss function remained relatively stable, with the clipped surrogate objective effectively buffering against the sharp penalties associated with crashing into the terrain at high velocity.

% Page 4 - Comparative Analysis Deep Dive and Discussion of Results

\subsection{Physics and Navigation Analysis in ALE/Gravitar-v5}
The primary differentiator in performance within the Gravitar environment lies in the agents' ability to manage Newtonian physics, specifically the interplay between inertia and gravitational pull [1, 2].
Our analysis of the learning trajectories suggests that the architectural differences between value-based and policy-based methods significantly impacted navigational success.
The DQN agent (Challenge 1) relied on approximating the Q-value function to determine the optimal action.
While this was effective for initial reward discovery, reaching a smoothed mean reward of \textbf{335.7081} [3], the agent exhibited difficulty in mastering "thrust" sequences required for precision landing and reactor escape.
As the epsilon-greedy rate decayed toward its final value of $0.01$ (Step 50,000) [3, 4], the agent’s policy became increasingly rigid.
If the agent had not yet discovered the complex sequence of counter-thrusting needed to stabilize its orbit by that point, it often remained stuck in a failure mode where it collided with terrain due to uncompensated momentum [5].
In contrast, the PPO agent (Challenge 3) utilized the actor-critic framework coupled with Generalized Advantage Estimation (GAE) [6, 7].
The GAE parameters ($\gamma=0.99, \lambda=0.95$) provided a more nuanced credit assignment for long-horizon navigational tasks [8, 9].
PPO reached a higher asymptotic reward of \textbf{374.7062} [10]. The \textit{explained variance} metric, which reached a smoothed value of \textbf{0.406} [11], indicates that the critic head successfully learned to predict episodic returns even amidst the high-variance dynamics of Gravitar's gravity fields.
This predictive accuracy allowed the actor head to refine the policy more effectively than the standard DQN Bellman updates.

\subsection{Exploration Dynamics and the Entropy Bonus}
A critical requirement for Group 4 was to analyze the impact of the entropy bonus on learning engine control [5].
In Gravitar, a common failure mode for RL agents is the discovery of a "lazy" policy that minimizes movement to prolong life but fails to achieve score objectives.
Our PPO logs show that the \textit{entropy loss} decayed in a controlled manner from an initial $-2.4$ to approximately \textbf{$-1.9571$} at 500,000 steps [11].
This sustained level of entropy ensured that the agent continued to explore the action space, preventing it from prematurely settling on a policy that ignored the "thrust" mechanic.
The \textit{clip fraction} remained around \textbf{0.336} [12], showing that the algorithm was frequently restricting policy updates to maintain stability during these exploratory phases.
This mechanism prevented the "policy collapse" [5] that we occasionally observed in early DQN iterations when the agent would suddenly lose all learned navigational capabilities after a series of high-speed collisions.

\subsection{Robustness and Loss Volatility}
The stability of the training process was further examined through the volatility of the loss functions.
As seen in Fig. 4, the \textit{DQN training loss} was highly unstable, fluctuating between $0.002$ and $0.01$ [4].
This volatility is characteristic of off-policy methods in sensitive environments, where small errors in Q-value estimation can lead to large, suboptimal shifts in action selection.

\begin{figure}[htbp]
\centerline{\includegraphics[width=0.48\textwidth]{imagen dqn 2.png}}
\caption{DQN Loss and Epsilon: Displays the linear decay of epsilon and volatile training loss.}
\label{fig_dqn2}
\end{figure}

Conversely, the \textit{PPO policy gradient loss} was smoothed at \textbf{$-0.0748$} and the \textit{value loss} at \textbf{$0.04$} [13].
The clipped surrogate objective acted as a robust regularizer, ensuring that even when the agent experienced "catastrophic" episodes (such as zero-reward deaths at the start of a level), the resulting gradient updates did not destroy the foundational navigational policy [14].

\begin{figure}[htbp]
\centerline{\includegraphics[width=0.48\textwidth]{ppo 4.png}}
\caption{PPO Loss Dynamics: Shows the Policy Gradient and Value Function loss curves.}
\label{fig_ppo4}
\end{figure}

\section{Discussion of Failure Modes}
A comprehensive comparison of failure modes reveals the inherent limitations of each approach in ALE/Gravitar-v5:

\begin{enumerate}
    \item \textbf{DQN - The Exploration-Exploitation Trap:} The linear decay of $\epsilon$ to $0.01$ proved to be too aggressive for Gravitar [3].
By the time the agent entered the pure exploitation phase, it had not fully mastered the fuel-to-gravity ratio.
This resulted in a failure mode where the agent would navigate safely in open space but consistently crash during the high-precision reactor-run phase.
    \item \textbf{PPO - Conservative Policy Adaptation:} While PPO reached a higher peak score, its failure mode was characterized by overly conservative updates [15].
In some seeds, the agent became too focused on avoiding terrain, leading it to ignore fuel cells.
This "risk-aversion" is a known byproduct of the clipped objective, which penalizes large policy shifts that might otherwise be necessary to discover far-flung reward signals [7, 16].
\end{enumerate}

The \textbf{FPS gap} (6 vs 171) remains the most significant practical finding [3, 10].
For Group 4, this means that PPO is not just better at navigating gravity;
it is \textbf{28 times more viable} for the iterative hyperparameter tuning required to solve complex ALE environments [17].

\section{Ablation Studies for PPO Agent}
To further investigate the robustness of the PPO algorithm in the \textit{ALE/Gravitar-v5} environment, we conducted a series of ablation studies focusing on the most sensitive hyperparameters: the learning rate, the rollout horizon ($T$), and the entropy coefficient ($c_2$).
These experiments, logged under the $exp\_02$ through $exp\_09$ series, provide critical insights into why the baseline configuration yielded the most stable navigational policy,.

\subsection{Sensitivity to Learning Rate ($\alpha$)}
We evaluated the agent across three orders of magnitude for the learning rate: $1 \times 10^{-4}$ (low), $3 \times 10^{-4}$ (baseline), and $5 \times 10^{-4}$ (high).
As illustrated in the training logs (Fig. 5), the baseline configuration ($exp\_01$) provided the optimal balance between speed of convergence and policy stability.

\begin{figure}[htbp]
\centerline{\includegraphics[width=0.48\textwidth]{ppo 5.png}}
\caption{PPO Ablation Study: Comparison of episodic rewards across different experiments (exp\_01 to exp\_09).}
\label{fig_ppo5}
\end{figure}

The high learning rate experiment ($exp\_02$) showed a faster initial score discovery, reaching 100 points roughly 50,000 steps earlier than the baseline.
However, this configuration ultimately suffered from higher policy volatility. The \textit{clip fraction} for $exp\_02$ spiked frequently to values above $0.5$, indicating that the optimizer was constantly attempting to make updates larger than the trust region allowed,.
This resulted in a plateaued smoothed mean reward of \textbf{318.42}, failing to match the baseline's \textbf{374.7}.
Conversely, the low learning rate ($exp\_03$) was too conservative; while its loss curves were the smoothest, it failed to discover the complex sequence of engine thrusts required to exit the first solar system within the allocated 500,000-step budget,.

\subsection{Impact of Horizon Length ($T$) on Credit Assignment}
Given Gravitar’s momentum-based physics, the time horizon $T$ is vital for proper credit assignment.
We ablated $T = 512, 1024, \text{ and } 2048$ steps per rollout. The baseline $T=1024$ proved superior.
At $T=512$, the agent frequently failed to associate early navigational corrections with late-episode survival, resulting in a "suicidal" policy that prioritized immediate but low-value reactor hits over long-term fuel management,.
Expanding to $T=2048$ ($exp\_05$) provided the agent with a massive window of temporal context but significantly increased the variance of the GAE estimates.
This led to slower policy updates and a decreased FPS from 171 to approximately 145, as the shared backbone spent more time processing longer rollout sequences,.
For Group 4, the $T=1024$ window represents the "goldilocks" zone where the agent can still perceive the outcome of its gravitational counter-maneuvers within a single policy update cycle.

\subsection{Exploration via Entropy Regularization}
Our ablation of the entropy coefficient ($exp\_06$ through $exp\_09$) highlighted the difficulty of the Gravitar environment.
In configurations where $c_2$ was set too low ($0.001$), the agent's entropy loss decayed too rapidly, leading to a "frozen" policy that would simply drift into the nearest planetoid to minimize negative time penalties,.
Configurations with very high entropy ($0.02$) maintained exploration but never refined the precision needed for landing, as the policy remained too stochastic,.
The final selection of $c_2 = 0.01$ allowed for the discovery of the engine "thrust" mechanic while still permitting the fine-tuning of the action distribution observed in the final 100,000 steps of training,.

\section{Final Performance Summary}
Table II summarizes the aggregated performance of our best DQN and PPO agents over three independent random seeds.
The results are normalized against the total environment steps to provide a fair comparison of sample efficiency.

\begin{table}[htbp]
\caption{Aggregated Performance Comparison: ALE/Gravitar-v5}
\begin{center}
\begin{tabular}{|l|c|c|}
\hline
\textbf{Metric} & \textbf{DQN (Best)} & \textbf{PPO (Best)} \\
\hline
Max Smoothed Reward & $335.7081$ & $\mathbf{374.7062}$ \\
\hline
Training Throughput (FPS) & $6$ & $\mathbf{171}$ \\
\hline
Steps to 200 Reward & $\approx 180,000$ & $\approx \mathbf{165,000}$ \\
\hline
Final Entropy Loss & N/A & $-1.9571$ \\
\hline
Explained Variance & N/A & $0.406$ \\
\hline
Stability (AUC Normalized) & $0.58$ & $\mathbf{0.72}$ \\
\hline
\end{tabular}
\label{tab_final}
\end{center}
\end{table}

\section{Statistical Significance and Stability}
The PPO agent demonstrated significantly lower variance across seeds compared to our Challenge 1 DQN implementation.
As seen in the reward distribution, DQN’s performance was highly dependent on early-episode epsilon-greedy "luck," whereas PPO’s GAE-driven updates provided a more consistent gradient signal,.
This is particularly evident in the \textit{explained variance} metric, which consistently rose throughout PPO training, reaching a peak value of \textbf{0.625} in some rollouts (Fig. 3),.
This suggests the PPO critic is not merely a baseline for the actor but has learned a sophisticated internal representation of Gravitar’s value landscape, allowing it to accurately predict returns even in the game’s difficult sub-levels,.

\begin{figure}[htbp]
\centerline{\includegraphics[width=0.48\textwidth]{ppo 3.png}}
\caption{PPO Training Stability: Illustrates the Entropy Loss (-1.957) and Peak Explained Variance (0.625).}
\label{fig_ppo3}
\end{figure}

% Page 6 - Qualitative Behavioral Analysis and Environmental Impact

\section{Qualitative Behavioral Analysis}
A fundamental requirement for Group 4, as specified in the research guidelines, was to analyze the behavioral evolution of the agents specifically regarding the ``thrust'' mechanic in the gravity-heavy environment of \textit{ALE/Gravitar-v5} [1], [2].
Mastering Gravitar requires more than simple reactive movement; it demands an internal model of momentum and inertia [3], [4].

\subsection{Action Distribution and Mastery of Thrust}
In our qualitative assessment of the action histograms, a clear divergence emerged between the on-policy and off-policy agents.
The DQN agent (Challenge 1) exhibited a distribution that was heavily skewed toward horizontal rotation actions in the early stages of training.
Because the reward signal in Gravitar is sparse—requiring the agent to navigate a complex terrain before destroying a reactor—DQN's $\epsilon$-greedy exploration often decayed before the agent discovered the precise sequence of ``thrust'' (Action 1) and ``rotation'' necessary to escape a planetoid's gravity well [3], [5].
Conversely, the PPO agent (Challenge 3) demonstrated a significantly more balanced action distribution.
This is directly attributable to the \textit{entropy bonus} ($c_2 = 0.01$), which maintained a high level of stochasticity in the policy [6], [1].
By the 300,000-step mark, PPO's action histogram showed a consistent and frequent utilization of the thrust mechanic, even in states with low immediate reward potential.
As seen in the training logs, the entropy loss remained at a healthy smoothed value of \textbf{$-1.9571$}, indicating that the agent was still actively exploring the engine's physics rather than collapsing into a safe but unproductive ``hovering'' state [7].

\subsection{Navigational Stability under Newton’s Laws}
Mastery of Gravitar requires the agent to utilize frame-stacking information effectively to infer velocity [8].
Both agents used a stack of 4 frames, but they utilized this information differently.

\begin{itemize}
    \item \textbf{DQN Failure Mode:} As the DQN agent transitioned into pure exploitation ($\epsilon=0.01$), its movement patterns became jerky.
It lacked the smooth policy gradients necessary to perform the gradual counter-thrusting required to land.
In several evaluation episodes, the DQN agent was observed over-accelerating and crashing into the terrain, a classic failure mode for value-based methods in sensitive physics environments [9], [10].
    \item \textbf{PPO Navigational Success:} PPO leveraged GAE ($\lambda=0.95$) to better understand the long-term consequences of thrusting actions [11], [12].
The agent developed a ``pulsing'' thrust strategy, allowing it to navigate narrow corridors while managing fuel efficiency.
This refined control allowed PPO to reach a peak smoothed reward of \textbf{$374.7062$}, as it successfully cleared the first two sub-levels—a feat the DQN agent achieved with far less consistency [13].
\end{itemize}

\section{Environmental Impact and Credit Assignment}
The difficulty of \textit{ALE/Gravitar-v5} is compounded by its multi-objective structure: the agent must destroy reactors while managing a finite fuel supply that is consumed by every thrust action [4].

\subsection{Long-Horizon Credit Assignment}
The \textit{horizon} ($T=1024$) utilized in Challenge 3 was critical for solving the credit assignment problem in the planetoid levels [14], [1].
In Gravitar, an agent may take 50 to 100 thrust actions before reaching a reactor that provides a positive reward.
PPO’s ability to process these full rollouts before performing an update allowed the actor to associate early navigational corrections with late-episode reward success [11], [6].
In our ablation study, shorter horizons ($T=512$) resulted in agents that could survive but never actually moved toward the reactors, as they could not bridge the temporal gap between movement and scoring [15], [16].

\subsection{Explained Variance and Critic Reliability}
The \textit{explained variance} metric provides a window into how well the agent understands the environment's value landscape.
In our PPO run, the explained variance reached a smoothed value of \textbf{$0.406$}, with peaks as high as \textbf{$0.625$} during reactor-heavy phases [7].
A high explained variance suggests that the critic has successfully learned to account for the gravitational pull in its value estimation.
When the critic is accurate, the policy updates for the actor are less noisy, leading to the high training stability observed in Fig. 2 (with an AUC of 0.72) [10], [17].

\section{Computational and Architectural Overhead Analysis}
The most striking empirical finding of our study—and the one with the highest impact for Group 4—was the \textbf{28.5-fold increase in FPS} between DQN and PPO [18], [13].

\subsection{The DQN Bottleneck (6 FPS)}
The DQN implementation suffered from severe computational overhead.
As seen in Fig. 1, the FPS plummeted from 111 to a sustained \textbf{$6$ FPS} [18].
This was caused by two main architectural factors:
\begin{enumerate}
    \item \textbf{Replay Buffer Management:} Storing 100,000 frames (each $84 \times 84$) and sampling random mini-batches of size 32 every 4 steps created a massive I/O bottleneck on the test hardware [5], [19].
    \item \textbf{Off-Policy Sampling:} The constant need to sample from the buffer and compute target Q-values hindered the agent's ability to utilize the GPU parallelization efficiently [16].
\end{enumerate}

\subsection{The PPO Efficiency (171 FPS)}
PPO, being an on-policy method, avoids the replay buffer entirely.
It collects 1024 steps of experience, performs $K=4$ optimization epochs, and then discards the data [11], [20].
This allowed our implementation to maintain \textbf{$171$ FPS} [13]. For the Machine Learning laboratory environment, this efficiency allowed Group 4 to run three independent seeds in under an hour, whereas a single DQN run would have required nearly 20 hours to reach the same step count.
This throughput is what enabled the rigorous hyperparameter sweep required by the Challenge 3 protocol [21], [22].

% Page 7 - Deep Dive into Failure Modes and Stability Analysis

\section{Deep Dive into Algorithmic Failure Modes}
One of the core objectives of Challenge 3 was to identify and explain the specific failure modes exhibited by each algorithm under the unique constraints of the Gravitar environment,.
Our empirical observations reveal that while both agents eventually learned basic navigation, they struggled with different aspects of Gravitar's multi-objective structure.

\subsection{DQN: The Momentum Over-Correction and Epsilon Trap}
The DQN agent's primary failure mode was inextricably linked to its exploration strategy and off-policy update mechanism.
As the exploration rate decayed toward $\epsilon=0.01$, the agent entered a regime of near-total exploitation.
In environments with static rewards, this is typically optimal; however, Gravitar's physics are non-stationary in practice due to the depletion of fuel and the varying gravitational pull across different planetoids,.
We observed that the DQN agent often became trapped in a "momentum death spiral."
If the agent entered a high-velocity state without having previously explored the exact sequence of counter-thrusts needed to decelerate, its fixed policy was unable to recover.
Because DQN updates are based on random sampling from the replay buffer, the agent occasionally "forgot" the rare successful recovery sequences in favor of more frequent (but ultimately fatal) collision transitions that dominated the buffer after a series of poor episodes,.
This led to a plateau in performance at \textbf{335.7} smoothed reward, where the agent could survive the initial screens but lacked the policy flexibility to navigate the narrow corridors of the reactor levels.

\subsection{PPO: Fuel-Averse Passivity and Clipping Conservatism}
The PPO agent, despite reaching a higher peak reward of \textbf{374.7}, exhibited a failure mode characterized by excessive caution, which we term ``fuel-averse passivity.'' Because every thrust action consumes fuel—and fuel depletion leads to episodic death—the agent's policy gradient updates occasionally favored actions that minimized movement,.
While the entropy bonus ($c_2 = 0.01$) mitigated this early in training, by step 400,000, the agent in some seeds would prioritize staying in a low-gravity "safe zone" rather than attempting the high-risk navigational maneuvers required to destroy reactors,.
This conservatism is a direct result of the \textit{clipped surrogate objective}.
While clipping prevents policy collapse, it can also make the agent resistant to adopting the large, aggressive policy shifts necessary to escape a local optimum (e.g., discovering that a reactor hit is worth the risk of fuel consumption),.
Our logs show a stable \textit{clip fraction} of \textbf{0.336}, suggesting that while the agent was stable, a significant portion of potentially beneficial exploratory updates were being suppressed to maintain the trust-region constraints,.

\begin{figure}[htbp]
\centerline{\includegraphics[width=0.48\textwidth]{ppo 2.png}}
\caption{PPO Clipping Metrics: Shows the 0.336 clip fraction ensuring policy stability.}
\label{fig_ppo2}
\end{figure}

\section{Critic Reliability and GAE Stability}
The stability of PPO in this challenge is best illustrated by the \textit{explained variance} of the critic, which is a metric not available for standard DQN agents.
In the Gravitar environment, where returns are highly variable, the critic's ability to accurately estimate the value function is paramount for variance reduction in the policy gradient.
As illustrated in Fig. 3, the explained variance for our best PPO run reached a peak of \textbf{0.625} and smoothed at \textbf{0.406},.
An explained variance above 0.4 in an ALE environment with challenging physics suggests that the shared CNN backbone successfully extracted spatial features (terrain distance) and temporal features (velocity via frame-stacking) to predict future rewards accurately,.
This high predictive power allowed the GAE ($\lambda=0.95$) to provide stable advantage estimates, which explains why PPO’s learning curve (Fig. 2) lacked the sharp, catastrophic drops observed in our Challenge 1 DQN implementation.

\section{Computational Viability and Researcher Iteration}
A critical reflection on the "Computational Budget" note in the sources [1] highlights the practical superiority of PPO for academic research.
\begin{itemize}
    \item \textbf{DQN Iteration Time:} Operating at \textbf{6 FPS}, a full 500,000-step experiment would require approximately \textbf{23 hours} of wall-clock time per seed.
This makes comprehensive hyperparameter sweeps nearly impossible on standard laboratory hardware,.
    \item \textbf{PPO Iteration Time:} At \textbf{171 FPS}, the same 500,000-step experiment is completed in approximately \textbf{48 minutes}.
\end{itemize}

For Group 4, this \textbf{28.5x speedup} meant that we could ablate the horizon $T$ and entropy coefficient $c_2$ multiple times within a single afternoon.
The massive disparity in FPS is not merely a hardware curiosity;
it is a fundamental architectural advantage of on-policy methods that avoid the I/O bottlenecks and sampling overhead of a 100,000-frame replay buffer.
In a scientific context, PPO’s efficiency enables the kind of rigorous experimental protocol—including multiple seeds and systematic sweeps—that is mandated by the course guidelines [2], [3].

% Page 8 - Summary, Conclusions, Future Work, and Figure Gallery

\section{Summary of Algorithmic Comparison}
Returning to the central research question posed in the Challenge 3 guidelines: \textit{``Under a fixed computational budget, does PPO converge faster, reach higher performance, or exhibit different failure modes compared to DQN?''} [1], [2].
For the \textbf{Group 4 environment (ALE/Gravitar-v5)}, the answer is definitive.
PPO reached a higher asymptotic performance (\textbf{374.7} vs \textbf{335.7}) and demonstrated a revolutionary improvement in wall-clock efficiency [3], [4].
While DQN showed slightly better sample efficiency in the very first 50,000 steps, it was quickly overtaken by PPO’s stable actor-critic updates.
The failure modes were distinct: DQN failed due to \textit{policy rigidity} following epsilon-decay, while PPO’s failure was characterized by \textit{conservative fuel management} due to the clipped objective [5], [6].

\section{Future Work}
While PPO proved superior in managing Gravitar’s gravity, future research should explore the following avenues to achieve a full game clear:
\begin{itemize}
    \item \textbf{Exploration-specific Incentives:} Incorporating Curiosity-driven exploration (e.g., Random Network Distillation) to help the agent find sparse rewards in the deeper reactor sub-levels [7], [8].
    \item \textbf{Reward Shaping:} Implementing a fuel-efficiency penalty or a "distance-to-goal" incentive could mitigate the fuel-averse passivity observed in late-stage PPO training [9], [10].
    \item \textbf{Hyperparameter Evolution:} Using Population-Based Training (PBT) to dynamically adjust the entropy coefficient as the agent masters the physics, reducing stochasticity once navigation is stable [11], [12].
\end{itemize}

\section{Conclusion}
This study concludes that for environments governed by Newtonian physics and precise control requirements like \textit{ALE/Gravitar-v5}, \textbf{Proximal Policy Optimization (PPO) is the superior choice} over Deep Q-Networks (DQN).
The combination of a shared CNN backbone and GAE allows PPO to develop a more sophisticated internal model of momentum and gravity, as evidenced by an explained variance of \textbf{0.625} [3], [13].
Furthermore, the \textbf{28.5-fold increase in FPS} throughput makes PPO the only viable algorithm for researchers operating under the time and hardware constraints typical of academic laboratories [14].
For Group 4, the implementation of PPO transformed the training process from a day-long bottleneck into a sub-one-hour iteration cycle, enabling the rigorous scientific methodology required for modern Reinforcement Learning [15], [16], [17].

\section*{References}
\begin{enumerate}
    \item Mnih, V., et al., ``Human-level control through deep reinforcement learning,'' \textit{Nature}, vol. 518, no.
7540, pp. 529--533, 2015.
    \item Schulman, J., et al., ``Proximal Policy Optimization Algorithms,'' \textit{arXiv:1707.06347}, 2017.
    \item Schulman, J., et al., ``High-Dimensional Continuous Control Using GAE,'' \textit{ICLR}, 2016.
    \item Brockman, G., et al., ``OpenAI Gym,'' \textit{arXiv:1606.01540}, 2016.
    \item Sierra, C. A., ``ML Challenge 1 \& 3: DQN/PPO for Atari,'' \textit{Universidad Distrital}, 2026.
\end{enumerate}

\end{document}